\chapter{Discussion}
\label{chap:discussion}

The results in Chapter~\ref{chap:results} support a specific claim rather
than a universal one. The placement layer improves runtime when the CUDASTF
task graph exposes copy-saving alternatives to HEFT, and it falls back to
HEFT-like behavior when such alternatives are absent. The DVFS layer then
provides conservative energy control on top of the same placement framework.
Its effect is measurable but workload dependent, which is consistent with the
energy--performance trade-offs described in Chapter~\ref{ch:theory}.

\section{Interpreting the Placement Result}
\label{sec:discussion-placement}

The placement results support the main design choice of this thesis: HEFT
should remain the performance anchor, while locality-aware deviations should be
admitted only when there is sufficient copy-saving evidence. In the placement
performance experiment, the bandit scheduler improves geometric-mean latency
by \(9.989\%\) across the five reported benchmarks. The improvement is not
uniform. CG, Cholesky, and LU decomposition show the strongest gains, FDTD is a
moderate positive case, and MiniWeather remains near neutral.

This non-uniformity is expected for a residual placement policy. The scheduler
does not try to replace HEFT with a globally different placement strategy.
Instead, it asks whether a non-HEFT GPU can reduce predicted copy cost while
keeping the local finish-time penalty bounded. The locality-attribution
analysis confirms that this behavior is visible in the measurements. The
largest improvements coincide with large copy byte reductions: Cholesky reduces
copy bytes by \(38.220\%\), LU decomposition reduces copy bytes by \(23.905\%\),
and CG reduces copy bytes by \(7.692\%\). MiniWeather keeps the same measured
copy footprint and remains close to HEFT, which is consistent with the
interpretation that residual placement is useful only when the task graph
exposes reducible data movement.

The trace-based mechanism analysis in
Section~\ref{sec:results-trace-locality-mechanism} adds an important nuance to
this interpretation. Copy-byte reduction should not be read as a direct
decomposition of latency. It is a locality-volume metric, not a measured trace
of transfer time or critical-path delay. The no-deviation control shows why the
locality explanation remains plausible: when the same scheduler path is used but
non-HEFT residual placements are disabled, the result returns to the near-HEFT
regime. The task trace then explains where the larger latency effect comes from:
the accepted residual placements reduce idle and readiness gaps along the
repeated execution chain, even though the measured task work itself does not
become shorter.

The placement result should therefore be read as a workload-dependent locality
result rather than as a uniform scheduler replacement. This matters for a
runtime scheduler because a method that improves copy-sensitive workloads while
preserving near-neutral behavior on low-opportunity workloads is more useful
than one that pursues aggressive deviations everywhere.

\section{Interpreting the DVFS Result}
\label{sec:discussion-dvfs}

The WinDVFS result should be interpreted as a conservative
energy-control result rather than as a universal energy optimum. Across the
five final rows, WinDVFS reduces geometric-mean energy by \(5.318\%\) with a
geometric-mean latency overhead of \(0.479\%\). The geometric-mean EDP ratio is
0.951354. FDTD, CG, MiniWeather, LU decomposition, and Cholesky all reduce EDP
relative to their no-DVFS counterparts, but the size of the benefit remains
workload dependent.

This outcome is consistent with the proposal--commit design. The DVFS layer
does not directly optimize each task in isolation. It accumulates task-level
frequency intents into per-GPU windows and commits a lower frequency only when
guardrails indicate that the performance risk is acceptable. The DVFS clock and
power measurements are used as supporting diagnostics rather than as the main
claim. The final DVFS measurements show that the larger energy reductions are
consistent with lower average power in the selected configurations, but the
primary evidence remains the measured energy, latency, and EDP trade-off.
Cholesky is the clearest trade-off case: it reduces energy, but its latency
increase limits the EDP improvement. MiniWeather provides a contrasting case
where latency remains essentially unchanged and the energy reduction is
smaller. These differences show that the controller should be interpreted as a
guarded trade-off mechanism, not as a rule that every workload should be
down-clocked in the same way.

The DVFS evidence is therefore best viewed as showing that task-graph runtime
information can support safe, window-level frequency control. It does not show
that the current heuristic is globally optimal, nor that every benchmark should
be down-clocked. The measured benefit is modest, but it is obtained while
preserving the separation between placement and DVFS and while keeping the
aggregate latency cost small.

\section{Answering the Research Questions}
\label{sec:discussion-research-questions}

The first research question asks how a CUDASTF runtime scheduler can preserve
the performance advantages of HEFT while exploiting data locality. The answer
in this thesis is a HEFT-relative residual gate. HEFT provides the baseline GPU
choice, and the scheduler deviates only when the predicted copy saving is large
enough and the finish-time penalty is bounded. The placement performance and
locality results show that this design improves latency on copy-sensitive
workloads while remaining close to HEFT on neutral workloads.

The second research question asks whether observed latency changes can be
attributed to concrete locality behavior. The locality-attribution experiment
provides the strongest evidence for this question. The benchmarks with the
largest latency improvements also show substantial copy-byte reductions.
Conversely, the near-neutral benchmark shows no measured copy-byte reduction.
This alignment does not prove that copy traffic explains every millisecond of
runtime variation, but it gives a clear data-movement account of the main
placement improvements.

The third research question asks how GPU DVFS decisions can be integrated into
the task scheduler without unstable per-task control. The proposal--commit
layer addresses this by separating task-level frequency intent from actual
frequency commitment. Tasks vote for high, mid, or low frequency, but the
committer changes the device-level committed level only at window boundaries
and only after guardrail checks.

The fourth research question asks whether the two-level policy can improve
energy while keeping runtime degradation bounded. The final DVFS result gives
a qualified positive answer. Energy decreases by \(5.318\%\) geometric mean,
while latency increases by only \(0.479\%\) geometric mean under the fixed
Bandit placement layer. The EDP ratio improves to 0.951354. The result is
workload dependent, so the correct conclusion is that the runtime can extract
energy savings on the reported benchmarks without large aggregate slowdown, not
that DVFS is always beneficial.

The fifth research question concerns overhead. The overhead diagnostics show
that the dynamic Bandit scheduler and the static-replay diagnostic differ by
only \(-0.766\%\) to \(+1.439\%\) in latency across the five reported
benchmarks, with a geometric-mean static-replay delta of \(+0.169\%\). This
supports the interpretation that online Bandit decision overhead is small at
benchmark scale. The HEFT replay row is shown only as a trace-replay sanity
check; the main overhead interpretation comes from the Bandit replay row.
Static replay is therefore treated as diagnostic evidence rather than as an
exact zero-overhead scheduler.

\section{Practical Implications}
\label{sec:discussion-practical-implications}

The main practical implication is that placement and DVFS should be separated
in an online multi-GPU runtime. Placement has an immediate effect on data
movement and queueing; DVFS affects power, frequency transitions, and possible
slowdown. Combining both into one large online action space would make reward
attribution difficult and would increase the risk of unstable behavior. The
layered design used here keeps each decision interpretable: placement can be
explained relative to HEFT and copy savings, while DVFS can be explained
through proposal shares, committed levels, clock exposure, and guardrail
events.

The results also suggest that non-regression behavior is as important as peak
improvement. A scheduler that improves only a subset of benchmarks is still
useful if it reliably avoids harming low-opportunity workloads. In this thesis,
MiniWeather is important precisely because it is a near-neutral control. It
shows that the scheduler does not need to force placement changes where the
measured locality opportunity is small.

\section{Limitations}
\label{sec:discussion-limitations}

The main limitation is experimental scope. The final measurements are taken on
a single-node eight-GPU configuration. The results may change on other GPU
models, interconnect topologies, CUDA driver versions, or problem sizes. In
particular, data-locality benefits depend on the relative cost of computation,
copy traffic, and queueing. A topology with different peer bandwidth or copy
engine behavior may change the magnitude of the placement gains.

The workload set is also limited. The five reported benchmarks cover several
useful cases, including strong locality cases, a near-neutral placement
control, and DVFS-positive rows, but they do not exhaust the range of CUDASTF
applications. The FDTD results also show that benchmark-level noise can affect
the apparent magnitude of small effects. This is why the thesis treats FDTD as
a moderate placement case rather than as the primary proof point.

The placement model uses a lightweight residency and copy-cost approximation.
This is appropriate for an online scheduler, but it cannot perfectly model
peer-copy overlap, copy-engine contention, topology-dependent bandwidth, or
all interactions with launch overheads. Similarly, the criticality
proxy used by placement and DVFS is local rather than a full dynamic critical
path computation. These simplifications keep the method practical, but they
also limit how precisely the scheduler can predict future global makespan.

Finally, the DVFS evaluation uses energy measured over the benchmark execution
window as the primary energy metric. This is the most relevant metric for the
frequency-control effect studied in this thesis, but it is not a full-system
energy study. A complete deployment evaluation would need longer runs, more
hardware platforms, full-system power measurements, and a more detailed
analysis of transition costs and background runtime overhead.
